\section{Introduction}
Matrix transposition exposes a distinction between identifying the required
permutation and paying to execute it. The distinction is especially sharp on
a machine with a fixed number of tapes: a short description of a permutation
does not by itself supply an efficient tape computation.

Our earlier papers established a native transposition model and studied local
coding and tape geometry within that model
\cite{FirstPaper2026,SecondPaper2026}. The original formalized results separate two questions.
First, can a linear-time bit transposer supply the word transposer used in the
canonical problem? Second, which restricted compilers could run in the opposite
direction? The constructive analysis uses a constant-length factorization of
binary address permutations. The negative analysis charges both exceptional
word operations and information carried across scan boundaries.

The barriers are statements about explicitly defined compiler grammars. They
remain valid when public operation descriptions vary with the dimension and
when permitted scan transitions are nonlinear. They are not normal-form
theorems for arbitrary multitape computations. In particular, they do not
resolve whether a more general word-to-bit reduction exists.

Harvey and van der Hoeven study reductions from transposition to integer
multiplication in the multitape model \cite{HarveyVanDerHoeven2025}. Section~7.6 of their 2025 version identifies the difficulty of transferring complexity bounds between
bit and multi-bit entries, and gives a bit-to-word simulation for oblivious
machines. Our reduction concerns the canonical logarithmic width
and the linear-time regime, while allowing an arbitrary adaptive bit machine.
It does not assert their more general comparison across all widths and time
bounds. The converse compiler barriers here also do not strengthen their
multiplication-based transposition algorithms: an unrestricted transposition lower bound would
still be required. These distinctions are necessary when interpreting progress
on the original multiplication problem.

The underlying methods have established precedents. Address-bit
factorizations occur in out-of-core transposition, for example in
\cite[Section~4]{KrishnamoorthyTranspose}; history-and-replay cleanup is due
to Bennett~\cite{Bennett1973}; and our shared-helper bound is an application
of the acyclic index-coding bound~\cite{BBJK2011}. The constructive result
combines its specific factorization with paid native execution to remove an
obliviousness assumption in the canonical linear-time setting.

The negative results concern different restrictions from earlier lower
bounds. Kirchherr's conservative-tape result treats entries as indivisible
pebbles~\cite{Kirchherr1988}. Adler et al. connect oblivious I/O transposition
with a network-coding conjecture~\cite[Section~6]{Adler2006}. Our scan bound
allows nonlinear mixing and requires no coding conjecture, but restricts
geometry, synchronization and state. It does not strengthen either result in its original model. Its contribution is a common-scale bound
and a fully charged helper derived from the concrete scan semantics.

The constructive extension in \cref{sec:phase} separates a further pair of
tasks. For dyadic matrices, a shrinking column recurrence computes a cyclic
phase and the least period in linear time, with all native tape costs paid.
Applying the resulting phases to the entire matrix is equivalent to canonical
word transposition. Thus phase discovery can be removed as an obstacle for
this particular normalization scheme, while the cost of producing the
normalized payload retains the full transposition question. These additional
results have complete written proofs, including the native schedules, but
have not been formalized in Lean. The original reduction and compiler
barriers retain their Lean proofs.
